% ============================================================================
% NGUỒN DUY NHẤT các gói & thư viện cần để DỰNG NỘI DUNG CÂU HỎI.
%
% Mọi đường biên dịch của app đều `\input` file này:
%   - xem trước câu hỏi / hình TikZ rời / Thư viện TikZ  → qua preamble.tex
%   - xuất đề PDF + phiếu trả lời                        → qua khaibaochung.tex
%
% VÌ SAO CÓ FILE NÀY: trước đây mỗi đường tự khai một danh sách gói riêng (preamble.tex,
% hằng TIKZ_LIBRARY_PACKAGES trong TypeScript, khaibaochung.tex). Ba danh sách lệch nhau
% sinh ra triệu chứng rất khó lần: CÙNG một câu render được ở trang này mà hỏng ở trang kia.
% Đo trên bank ngày 21/07/2026: 445 câu (~1%) dính đúng lỗi đó.
%
% QUY TẮC: thêm gói/thư viện phục vụ NỘI DUNG câu hỏi thì thêm VÀO ĐÂY, đừng thêm vào
% preamble.tex hay khaibaochung.tex. Script `npm run check:preamble` sẽ báo đỏ nếu hai file
% kia tự khai `\usetikzlibrary`/`\usepgfplotslibrary` ngoài file này.
% Sửa file này PHẢI tăng PREAMBLE_VERSION trong src/main/latex/preamble.ts (nếu không,
% câu đã cache vẫn trả về bản render cũ và bản vá trông như vô tác dụng).
%
% KHÔNG thuộc phạm vi file này (cố ý để riêng mỗi đường): documentclass, engine biên dịch,
% font, hyperref, bố cục đề (titlesec/fancyhdr/counters). Xem docs/latex-format.md mục 4.
% ============================================================================

% --- Bộ vẽ hình ---
\usepackage{tikz-3dplot}          % hình không gian 3D (\tdplotsetmaincoords, hình chóp/lăng trụ)
\usepackage{tikz}
\usepackage{tkz-euclide}
\usepackage{tkz-tab}              % bảng biến thiên / xét dấu (các câu image_type=center)
\usepackage{tabvar}               % biến thể khác của bảng biến thiên (một số câu cũ)
\usepackage{pgf,pgfplots}

% Hợp của cả ba danh sách cũ. Số câu trong bank cần từng thư viện (đo 21/07/2026) ghi kèm
% những thư viện TỪNG THIẾU ở preamble.tex — đó chính là các câu hỏng khi xem trước:
%   intersections           190 câu (`name path=`)
%   decorations.markings     87 câu (`decoration={markings,…}`)
%   arrows                   60 câu (mũi tên cũ `->|`, `stealth'` — arrows.meta KHÔNG có)
%   positioning               3 câu (`right=of …`)
%   through                   2 câu
% `hobby`, `decorations.pathreplacing`: 0 câu trong bank nhưng người dùng có thể gõ tay ở
% Thư viện TikZ / Trợ lí TeX nên vẫn giữ.
% CỐ Ý KHÔNG có `snakes`: thư viện đã khai tử và xung khắc với `decorations.*`. 13 câu dùng
% `decoration={snake,…}` là cú pháp MỚI do `decorations.pathmorphing` cung cấp; câu duy nhất
% còn cú pháp cũ `snake=coil` (2D4H2-6) đã xóa mềm ngày 21/07/2026.
\usetikzlibrary{%
  angles,arrows,arrows.meta,backgrounds,calc,%
  decorations.markings,decorations.pathmorphing,decorations.pathreplacing,%
  hobby,intersections,patterns,patterns.meta,positioning,quotes,shadows,%
  shapes.geometric,shapes.symbols,through%
}
% compat: để 1.17 thay vì 1.18 để tương thích với máy người dùng có PGFPlots cũ (TeX Live
% 2020/2021, MiKTeX chưa cập nhật). PGFPlots 1.18+ vẫn chạy tốt với compat=1.17.
% Đã đối chiếu 23 câu dùng pgfplots — không đổi kết quả đáng kể.
\pgfplotsset{compat=1.17}
\usepgfplotslibrary{fillbetween,polar}

% --- Gói nội dung khác mà câu hỏi trong bank có dùng ---
% Trước đây chỉ khaibaochung.tex có ⇒ các câu này hỏng khi xem trước.
\usepackage{multirow,makecell,diagbox}  % bảng ghép ô (7+3+3 câu)
\usepackage{fontawesome,pifont}         % \faXXX (56 câu), \ding (2 câu)
\usepackage{yhmath}                     % \wideparen — cung tròn (8 câu)
\usepackage{esvect}                     % \vv (128 câu gõ tay)
\usepackage{twemojis}                   % \twemoji (5 câu); .sty nằm cùng thư mục này

% CỐ Ý KHÔNG chép `\def\vec{\vv}` / `\def\overrightarrow{\vv}` của khaibaochung.tex sang đây:
% đã thử và gây "Undefined control sequence" hàng loạt trên 5.721 câu dùng \overrightarrow
% chuẩn (xem ghi chú trong preamble.tex). Nạp gói esvect là đủ cho các câu gõ thẳng \vv.

% --- Hộp ghi chú (tcolorbox) ---
% Đo được +4 ms mỗi lượt biên dịch nên nạp cho mọi đường, đổi lại `\begin{dl}`, `\begin{luuy}`…
% hiện GIỐNG NHAU ở xem trước và ở đề in ra.
\usepackage{tcolorbox}
\tcbuselibrary{theorems,breakable,skins}

\newtcolorbox{khung}{
  before upper=\raggedright,
  colback=white,
  boxrule=0.5pt,
  colframe=gray!50!black,
  fonttitle=\bfseries,
  breakable,
  arc=1mm,
  left=2pt, right=2pt, top=2pt, bottom=2pt
}

% Hộp nền màu, viền trái — mỗi loại ghi chú một màu.
\newenvironment{boxdn}
{\begin{tcolorbox}[enhanced jigsaw,breakable,pad at break*=1mm,colback=red!10,boxrule=0pt,
    frame hidden,opacityback=0.5,left=3mm,right=0pt,bottom=2pt,top=2pt,before skip=2mm,
    after skip=2mm,arc=0pt,borderline west={1mm}{0cm}{red!30}]}
{\end{tcolorbox}}
\newenvironment{boxdl}
{\begin{tcolorbox}[enhanced jigsaw,breakable,pad at break*=1mm,colback=cyan!10,boxrule=0pt,
    frame hidden,opacityback=0.5,left=3mm,right=0pt,bottom=2pt,top=2pt,before skip=2mm,
    after skip=2mm,arc=0pt,borderline west={1mm}{0cm}{cyan!30}]}
{\end{tcolorbox}}
\newenvironment{boxchuy}
{\begin{tcolorbox}[enhanced jigsaw,breakable,pad at break*=1mm,colback=green!10,boxrule=0pt,
    frame hidden,opacityback=0.5,left=3mm,right=0pt,bottom=2pt,top=2pt,before skip=2mm,
    after skip=2mm,arc=0pt,borderline west={1mm}{0cm}{green!30}]}
{\end{tcolorbox}}
\newenvironment{boxluuy}
{\begin{tcolorbox}[enhanced jigsaw,breakable,pad at break*=1mm,colback=purple!10,boxrule=0pt,
    frame hidden,opacityback=0.5,left=3mm,right=0pt,bottom=2pt,top=2pt,before skip=2mm,
    after skip=2mm,arc=0pt,borderline west={1mm}{0cm}{purple!30}]}
{\end{tcolorbox}}
\newenvironment{boxpp}
{\begin{tcolorbox}[enhanced jigsaw,breakable,pad at break*=1mm,colback=orange!10,boxrule=0pt,
    frame hidden,opacityback=0.5,left=3mm,right=0pt,bottom=2pt,top=2pt,before skip=2mm,
    after skip=2mm,arc=0pt,borderline west={1mm}{0cm}{orange!30}]}
{\end{tcolorbox}}
\newenvironment{boxtc}
{\begin{tcolorbox}[enhanced jigsaw,breakable,pad at break*=1mm,colback=pink!10,boxrule=0pt,
    frame hidden,opacityback=0.5,left=3mm,right=0pt,bottom=2pt,top=2pt,before skip=2mm,
    after skip=2mm,arc=0pt,borderline west={1mm}{0cm}{pink!30}]}
{\end{tcolorbox}}
% Hộp LÝ THUYẾT — dùng cho tính năng "OCR tài liệu" (gõ lại nguyên khối lý thuyết của sách/tài
% liệu). Cùng khuôn với các hộp trên; `breakable` là bắt buộc vì một khối lý thuyết thường dài hơn
% một trang. Nhãn nhận THAM SỐ TÙY CHỌN để giữ đúng tên mục của tài liệu gốc:
%   \begin{lt} …            → nhãn "Lý thuyết."
%   \begin{lt}[Hàm số bậc hai] … → nhãn "Hàm số bậc hai."
\newenvironment{boxlt}
{\begin{tcolorbox}[enhanced jigsaw,breakable,pad at break*=1mm,colback=teal!10,boxrule=0pt,
    frame hidden,opacityback=0.5,left=3mm,right=0pt,bottom=2pt,top=2pt,before skip=2mm,
    after skip=2mm,arc=0pt,borderline west={1mm}{0cm}{teal!30}]}
{\end{tcolorbox}}

\newenvironment{dn}{\begin{boxdn}\textcolor{red}{\textbf{Định nghĩa. }}}{\end{boxdn}}
\newenvironment{dl}{\begin{boxdl}\textcolor{blue}{\textbf{Định lý. }}}{\end{boxdl}}
\newenvironment{luuy}{\begin{boxluuy}\textcolor{purple}{\textbf{Lưu ý. }}}{\end{boxluuy}}
\newenvironment{nx}{\begin{khung}\textcolor{purple}{\textbf{Nhận xét. }}}{\end{khung}}
\newenvironment{pp}{\begin{boxpp}\textcolor{orange!70!black}{\textbf{Phương pháp. }}}{\end{boxpp}}
\newenvironment{chuy}{\begin{boxchuy}\textcolor{green!50!black}{\textbf{Chú ý. }}}{\end{boxchuy}}
\newenvironment{note}{\begin{khung}\textcolor{cyan}{\textbf{Ghi chú. }}}{\end{khung}}
\newenvironment{tc}{\begin{boxtc}\textcolor{pink!50!black}{\textbf{Tính chất. }}}{\end{boxtc}}
\newenvironment{lt}[1][Lý thuyết]{\begin{boxlt}\textcolor{teal!60!black}{\textbf{#1. }}}{\end{boxlt}}

% --- Macro tự chế của tác giả gốc (không thuộc ex_test.sty) ---
% Hệ/tuyển phương trình — không phải lệnh LaTeX chuẩn, nhiều câu trong bank dùng.
\newcommand{\hoac}[1]{\left[\begin{aligned}#1\end{aligned}\right.}
\newcommand{\heva}[1]{\left\{\begin{aligned}#1\end{aligned}\right.}
